\documentclass[11pt]{article}
\usepackage[utf8]{inputenc}
\usepackage[T1]{fontenc}
\usepackage{geometry}
\usepackage{enumitem}
\usepackage{amsmath}
\usepackage{xcolor}

\geometry{margin=1in}

\title{Week 5 Quiz: Answer Key}
\author{Prompt Engineering Course}
\date{}

\begin{document}

\maketitle

\section*{Answer Key - Multiple Choice (10 points each)}

\begin{enumerate}

\item \textbf{B. Field, tenor, mode}
\begin{itemize}
    \item Field = topic domain/technical complexity
    \item Tenor = relationship/power distance
    \item Mode = spoken/written, planned/spontaneous
\end{itemize}

\item \textbf{B. The relationship and power distance between communicators}
\begin{itemize}
    \item Tenor manages social relationships: authority, peer, subordinate
    \item Controls solidarity markers and formality level
\end{itemize}

\item \textbf{B. Sentence length, clause complexity, and passive voice percentage}
\begin{itemize}
    \item Syntax signature captures structural patterns
    \item Measurable metrics like avg.\ sentence length (12--18 words)
    \item Passive voice percentage ($<$10\%)
\end{itemize}

\item \textbf{B. AI learns boundaries better from contrasts than from positive examples alone}
\begin{itemize}
    \item Contrasts help define what to avoid, not just what to do
    \item Strong sample shows target; counter-example shows boundary
\end{itemize}

\item \textbf{B. A specification that defines exactly how output should be formatted for downstream processing}
\begin{itemize}
    \item Ensures outputs are parse-ready (JSON, CSV, XML)
    \item Enables automation and integration
    \item Includes validation hooks
\end{itemize}

\item \textbf{C. ``Use short sentences, action verbs, and explicit deadlines''}
\begin{itemize}
    \item Urgent tone = brevity + action + time pressure
    \item Opposite: long sentences, passive voice = slower, more neutral
\end{itemize}

\item \textbf{B. Use layered prompts with priority matrices}
\begin{itemize}
    \item Acknowledge different stakeholder needs
    \item Establish hierarchy when conflicts arise
    \item Create modular outputs if needed
\end{itemize}

\item \textbf{B. To automatically validate tone consistency using metrics like sentiment analysis and readability scores}
\begin{itemize}
    \item Self-labeling: AI reports its own tone choices
    \item Reference comparison: measure similarity to exemplars
    \item Automated quality control at scale
\end{itemize}

\item \textbf{C. Multi-dimensional specs provide precise, reproducible style control}
\begin{itemize}
    \item Vague instructions like ``be professional'' are subjective and inconsistent
    \item Multi-dimensional specs (field + tenor + mode) are measurable and reproducible
    \item Enables consistent AI outputs across multiple generations
\end{itemize}

\item \textbf{C. Whether communication is spoken/written and planned/spontaneous}
\begin{itemize}
    \item Mode = delivery channel and planning level
    \item Interactive vs.\ monologic
    \item Written-planned vs.\ spoken-spontaneous
    \item Affects sentence structure and completeness
\end{itemize}

\end{enumerate}

\vspace{1cm}
\noindent \textbf{Total Points: /100}

\end{document}
